\documentclass[aspectratio=169,11pt]{beamer}

% Packages
\usepackage[utf8]{inputenc}
\usepackage[T1]{fontenc}
\usepackage[french]{babel}
\usepackage{graphicx}
\usepackage{listings}
\usepackage{xcolor}
\usepackage{tikz}
\usepackage{fontawesome5}

% Thème
\usetheme{Madrid}
\usecolortheme{default}

% Couleurs personnalisées
\definecolor{primaryblue}{RGB}{41,128,185}
\definecolor{secondaryblue}{RGB}{52,152,219}
\definecolor{accentorange}{RGB}{231,76,60}
\definecolor{successgreen}{RGB}{39,174,96}

\setbeamercolor{structure}{fg=primaryblue}
\setbeamercolor{title}{fg=white,bg=primaryblue}
\setbeamercolor{frametitle}{fg=white,bg=primaryblue}

% Informations
\title[Robot Hotel]{Cahier des Charges}
\subtitle{Projet Robot Hotel - Système Intelligent d'Accueil et de Gestion Hôtelière}
\author{Équipe de Développement}
\institute{Université d'Avignon}
\date{\today}

% Logo (optionnel - à adapter selon votre logo)
% \logo{\includegraphics[height=0.8cm]{logo.png}}

\begin{document}

% Page de titre
\frame{\titlepage}

% Table des matières
\begin{frame}{Plan de la Présentation}
\tableofcontents
\end{frame}

% ============================================================================
% SECTION 1: CONTEXTE ET OBJECTIFS
% ============================================================================

\section{Contexte et Objectifs}

\begin{frame}{Contexte du Projet}
\begin{block}{Vision}
Système intelligent d'accueil et de gestion hôtelière basé sur robot humanoïde Pepper
\end{block}

\begin{columns}
\column{0.5\textwidth}
\textbf{Objectifs principaux :}
\begin{itemize}
    \item Améliorer l'expérience d'accueil
    \item Automatiser les tâches de réception
    \item Service d'information 24/7
    \item Gestion automatisée des réservations
\end{itemize}

\column{0.5\textwidth}
\textbf{Besoins identifiés :}
\begin{itemize}
    \item Personnalisation de l'accueil
    \item Accessibilité multilingue
    \item Efficacité opérationnelle
    \item Intégration technologique
\end{itemize}
\end{columns}
\end{frame}

\begin{frame}{Objectifs Fonctionnels}
\begin{block}{Fonctionnalités clés}
\begin{itemize}
    \item[\textcolor{primaryblue}{\faUser}] \textbf{Reconnaissance faciale} pour identification clients
    \item[\textcolor{primaryblue}{\faComments}] \textbf{Chatbot intelligent} (Llama 3.1) avec RAG
    \item[\textcolor{primaryblue}{\faMicrophone}] \textbf{Transcription audio} (Whisper)
    \item[\textcolor{primaryblue}{\faDatabase}] \textbf{API REST complète} (Flask)
    \item[\textcolor{primaryblue}{\faDesktop}] \textbf{Interface admin} (Streamlit)
    \item[\textcolor{primaryblue}{\faBed}] \textbf{Gestion réservations} automatisée
\end{itemize}
\end{block}

\begin{alertblock}{Objectifs techniques}
Architecture modulaire, performances temps réel, compatibilité Pepper, documentation API
\end{alertblock}
\end{frame}

% ============================================================================
% SECTION 2: PÉRIMÈTRE
% ============================================================================

\section{Périmètre du Projet}

\begin{frame}{Fonctionnalités Principales}
\begin{columns}
\column{0.5\textwidth}
\textbf{\faUser\ Reconnaissance Faciale}
\begin{itemize}
    \item Enregistrement personnes
    \item Détection temps réel
    \item Streaming vidéo
    \item Gestion base de données
\end{itemize}

\vspace{0.5cm}

\textbf{\faComments\ Chat Intelligent}
\begin{itemize}
    \item Chatbot Llama 3.1
    \item Contexte RAG
    \item Connaissance temporelle
    \item Support multilingue
\end{itemize}

\column{0.5\textwidth}
\textbf{\faMicrophone\ Audio}
\begin{itemize}
    \item Transcription Whisper
    \item Enregistrement vocal
    \item Détection silence
    \item Extraction informations
\end{itemize}

\vspace{0.5cm}

\textbf{\faBed\ Réservations}
\begin{itemize}
    \item Création depuis audio
    \item Extraction intelligente
    \item QR codes automatiques
    \item Vérification QR codes
\end{itemize}
\end{columns}
\end{frame}

% ============================================================================
% SECTION 3: TECHNOLOGIES
% ============================================================================

\section{Technologies Utilisées}

\begin{frame}{Stack Technologique}
\begin{columns}
\column{0.5\textwidth}
\textbf{Frameworks}
\begin{itemize}
    \item Flask (API REST)
    \item FastAPI (Réservations)
    \item Streamlit (Interface)
    \item Uvicorn (Serveur ASGI)
\end{itemize}

\vspace{0.5cm}

\textbf{IA / ML}
\begin{itemize}
    \item Llama 3.1-8B-Instruct
    \item Whisper Medium
    \item YOLOv8
    \item face\_recognition
    \item OpenCV
\end{itemize}

\column{0.5\textwidth}
\textbf{RAG}
\begin{itemize}
    \item LlamaIndex
    \item HuggingFace Embeddings
    \item Vector Store
\end{itemize}

\vspace{0.5cm}

\textbf{Bases de Données}
\begin{itemize}
    \item Pickle (Visages)
    \item SQLite (Réservations)
    \item SQLAlchemy (ORM)
\end{itemize}

\vspace{0.5cm}

\textbf{Matériel}
\begin{itemize}
    \item Robot Pepper
    \item Webcam
    \item Microphone
\end{itemize}
\end{columns}
\end{frame}

% ============================================================================
% SECTION 4: FONCTIONNALITÉS IMPLÉMENTÉES
% ============================================================================

\section{Fonctionnalités Implémentées}

\begin{frame}{Architecture Multi-Services}
\begin{block}{Système composé de 3 services principaux}
\begin{enumerate}
    \item \textbf{API Flask} (port 5000) - Serveur principal avec Swagger
    \item \textbf{API FastAPI} (port 8000) - Gestion réservations
    \item \textbf{Interface Streamlit} (port 8501) - Administration
\end{enumerate}
\end{block}

\begin{block}{Automatisation}
Script Bash unifié (\texttt{start\_all.sh}) pour démarrer tous les services avec leurs environnements virtuels
\end{block}
\end{frame}

\begin{frame}{Reconnaissance Faciale}
\begin{columns}
\column{0.5\textwidth}
\textbf{Technologies}
\begin{itemize}
    \item YOLOv8 (détection)
    \item face\_recognition (reconnaissance)
    \item Base pickle
\end{itemize}

\vspace{0.5cm}

\textbf{Fonctionnalités}
\begin{itemize}
    \item Enregistrement via upload
    \item Détection temps réel
    \item Streaming vidéo
    \item Gestion complète BDD
\end{itemize}

\column{0.5\textwidth}
\begin{center}
\Large
\textcolor{primaryblue}{\faUser\ \faCheckCircle}
\end{center}

\textbf{Interface Admin}
\begin{itemize}
    \item Page Streamlit dédiée
    \item Preview images
    \item Liste et suppression
    \item Test en temps réel
\end{itemize}
\end{columns}
\end{frame}

\begin{frame}{Chat Intelligent avec RAG}
\begin{block}{Chatbot Llama 3.1-8B-Instruct}
\begin{itemize}
    \item Connexion Hugging Face
    \item Conversations contextuelles
    \item Réponses basées sur documentation hôtel
\end{itemize}
\end{block}

\begin{block}{Système RAG}
\begin{itemize}
    \item Indexation vectorielle (LlamaIndex)
    \item Embeddings HuggingFace
    \item Recherche sémantique
\end{itemize}
\end{block}

\begin{block}{Contexte Temporel}
Intégration automatique date/heure pour réponses temporelles précises
\end{block}
\end{frame}

\begin{frame}{Gestion de Réservations}
\begin{columns}
\column{0.5\textwidth}
\textbf{Création depuis Audio}
\begin{itemize}
    \item Transcription Whisper
    \item Extraction Llama 3.1
    \item Conversion dates relatives
\end{itemize}

\vspace{0.5cm}

\textbf{API FastAPI}
\begin{itemize}
    \item Création clients
    \item Création réservations
    \item Vérification QR codes
    \item Confirmation arrivée
\end{itemize}

\column{0.5\textwidth}
\textbf{QR Codes}
\begin{itemize}
    \item Génération automatique
    \item Stockage PNG
    \item Vérification scan
    \item Interface dédiée
\end{itemize}

\vspace{0.5cm}

\textbf{Notifications Email}
\begin{itemize}
    \item Envoi automatique
    \item Email HTML formaté
    \item QR code en pièce jointe
    \item Configuration SMTP
\end{itemize}
\end{columns}
\end{frame}

\begin{frame}{Interface Streamlit - 5 Onglets}
\begin{block}{Organisation}
\begin{enumerate}
    \item \textcolor{primaryblue}{\faUser} \textbf{Gestion Visages} - Enregistrement, liste, suppression
    \item \textcolor{primaryblue}{\faVideo} \textbf{Streaming Vidéo} - Contrôle détection faciale
    \item \textcolor{primaryblue}{\faComments} \textbf{Chat et Audio} - Chat, transcription, enregistrement
    \item \textcolor{primaryblue}{\faBed} \textbf{Réservations} - Création, vérification QR codes
    \item \textcolor{primaryblue}{\faEnvelope} \textbf{Configuration Email} - Paramètres SMTP
\end{enumerate}
\end{block}

\begin{block}{Fonctionnalités}
Test connexion API, gestion erreurs, feedback utilisateur, logs détaillés
\end{block}
\end{frame}

\begin{frame}{Intégration Robot Pepper}
\begin{columns}
\column{0.5\textwidth}
\textbf{Script Principal}
\begin{itemize}
    \item Interaction complète
    \item Enregistrement audio
    \item Transcription + Chat
    \item Réponse vocale
\end{itemize}

\vspace{0.5cm}

\textbf{Fonctionnalités}
\begin{itemize}
    \item Création réservations
    \item Scan QR codes
    \item Interface tablette
\end{itemize}

\column{0.5\textwidth}
\begin{center}
\Large
\textcolor{accentorange}{\faRobot\ \faCheckCircle}
\end{center}

\textbf{Workflow}
\begin{enumerate}
    \item Reconnaissance
    \item Interaction vocale
    \item Traitement IA
    \item Réponse robot
\end{enumerate}
\end{columns}
\end{frame}

% ============================================================================
% SECTION 5: SCÉNARIOS DE TESTS
% ============================================================================

\section{Scénarios de Tests}

\begin{frame}{Scénarios Principaux}
\begin{block}{Scénario 1 : Reconnaissance}
\begin{itemize}
    \item Enregistrement visage
    \item Détection robot/caméra
    \item Accueil personnalisé
    \item \textbf{Critère} : 80\% précision, < 3s
\end{itemize}
\end{block}

\begin{block}{Scénario 2 : Chat}
\begin{itemize}
    \item Questions services/équipements
    \item Questions temporelles
    \item Questions multilingues
    \item \textbf{Critère} : > 85\% pertinence, < 5s
\end{itemize}
\end{block}
\end{frame}

\begin{frame}{Scénarios Principaux (suite)}
\begin{block}{Scénario 3 : Réservation Vocale}
\begin{itemize}
    \item Enregistrement demande
    \item Transcription automatique
    \item Extraction informations
    \item Génération QR code
    \item \textbf{Critère} : > 90\% transcription, > 85\% extraction
\end{itemize}
\end{block}

\begin{block}{Scénario 4 : Administration}
\begin{itemize}
    \item Navigation interface
    \item Gestion visages
    \item Tests fonctionnalités
    \item \textbf{Critère} : Navigation intuitive, < 2min/opération
\end{itemize}
\end{block}
\end{frame}

% ============================================================================
% SECTION 6: CRITÈRES DE RÉUSSITE
% ============================================================================

\section{Critères de Réussite}

\begin{frame}{Performances Techniques}
\begin{block}{Indicateurs Clés}
\begin{itemize}
    \item \textbf{Reconnaissance faciale} : $\geq$ 85\% précision
    \item \textbf{Temps détection} : $\leq$ 3 secondes
    \item \textbf{Transcription} : $\geq$ 90\% (WER)
    \item \textbf{Réponse API} : $\leq$ 2s (95\% requêtes)
    \item \textbf{Disponibilité} : $\geq$ 95\%
\end{itemize}
\end{block}

\begin{block}{Performances Utilisateur}
\begin{itemize}
    \item Taux réussite tâches : $\geq$ 80\%
    \item Temps réservation : $\leq$ 2 minutes
    \item Temps question chatbot : $\leq$ 30 secondes
    \item Taux abandon : $\leq$ 15\%
\end{itemize}
\end{block}
\end{frame}

\begin{frame}{Satisfaction Utilisateur}
\begin{block}{Métriques}
\begin{itemize}
    \item \textbf{Score global} : $\geq$ 4.0/5.0
    \item \textbf{Net Promoter Score} : $\geq$ 50
    \item \textbf{Intention utilisation future} : $\geq$ 70\%
    \item \textbf{Recommandation} : $\geq$ 75\%
\end{itemize}
\end{block}

\begin{block}{Collecte de Retours}
\begin{itemize}
    \item Questionnaires structurés (pré/post-test)
    \item Entretiens semi-directifs
    \item Observation directe (vidéo)
    \item Métriques comportementales
\end{itemize}
\end{block}
\end{frame}

% ============================================================================
% SECTION 7: AMÉLIORATIONS FUTURES
% ============================================================================

\section{Améliorations Futures}

\begin{frame}{Court Terme (3-6 mois)}
\begin{block}{Priorités}
\begin{itemize}
    \item[\textcolor{successgreen}{\faCheckSquare}] Dashboard analytics
    \item[\textcolor{successgreen}{\faCheckSquare}] Authentification multi-utilisateurs
    \item[\textcolor{successgreen}{\faCheckSquare}] Amélioration précision reconnaissance
    \item[\textcolor{successgreen}{\faCheckSquare}] Templates emails personnalisables
    \item[\textcolor{successgreen}{\faCheckSquare}] Dockerisation complète
\end{itemize}
\end{block}
\end{frame}

\begin{frame}{Moyen Terme (6-12 mois)}
\begin{block}{Évolutions}
\begin{itemize}
    \item[\textcolor{secondaryblue}{\faMobile}] Application mobile clients
    \item[\textcolor{secondaryblue}{\faBrain}] Fine-tuning chatbot
    \item[\textcolor{secondaryblue}{\faCreditCard}] Système de paiement
    \item[\textcolor{secondaryblue}{\faSms}] Notifications SMS/WhatsApp
    \item[\textcolor{secondaryblue}{\faDatabase}] Migration PostgreSQL
    \item[\textcolor{secondaryblue}{\faChartLine}] Monitoring et alertes
\end{itemize}
\end{block}
\end{frame}

\begin{frame}{Long Terme (12+ mois)}
\begin{block}{Vision}
\begin{itemize}
    \item[\textcolor{accentorange}{\faRocket}] IA avancée (sentiment, recommandations)
    \item[\textcolor{accentorange}{\faUsers}] Support multi-robots
    \item[\textcolor{accentorange}{\faPlug}] Intégrations externes (PMS, CRM)
    \item[\textcolor{accentorange}{\faChartBar}] Analytics prédictifs
    \item[\textcolor{accentorange}{\faGlobe}] Extension autres secteurs
\end{itemize}
\end{block}
\end{frame}

% ============================================================================
% SECTION 8: CONTRAINTES ET RISQUES
% ============================================================================

\section{Contraintes et Risques}

\begin{frame}{Contraintes Techniques}
\begin{columns}
\column{0.5\textwidth}
\textbf{Matérielles}
\begin{itemize}
    \item Disponibilité Pepper
    \item Ressources computationnelles
    \item Stockage modèles ML
    \item Connexion réseau
\end{itemize}

\column{0.5\textwidth}
\textbf{Logicielles}
\begin{itemize}
    \item Python 3.12+
    \item Dépendances multiples
    \item Environnements virtuels
    \item Compatibilité OS
\end{itemize}
\end{columns}

\begin{alertblock}{Mesures}
Tests approfondis, documentation complète, logs détaillés, mode dégradé
\end{alertblock}
\end{frame}

\begin{frame}{Risques et Mitigation}
\begin{block}{Risques Techniques}
\begin{itemize}
    \item Performance modèles IA → Tests approfondis
    \item Intégration Pepper → Documentation complète
    \item Bugs critiques → Système logs détaillé
\end{itemize}
\end{block}

\begin{block}{Risques Utilisateur}
\begin{itemize}
    \item Acceptabilité → Tests utilisateurs précoces
    \item Complexité perçue → Design simple
    \item Fiabilité → Messages erreurs clairs
\end{itemize}
\end{block}
\end{frame}

% ============================================================================
% SECTION 9: PLANNING
% ============================================================================

\section{Planning Prévisionnel}

\begin{frame}{Phases du Projet - 24 semaines}
\begin{block}{Timeline}
\begin{enumerate}
    \item \textbf{Conception} (4 sem) - Architecture, technologies
    \item \textbf{Développement Core} (6 sem) - API, reconnaissance, chatbot
    \item \textbf{Fonctionnalités Avancées} (4 sem) - Réservations, optimisations
    \item \textbf{Tests et Intégration} (3 sem) - Tests techniques, Pepper
    \item \textbf{Tests Utilisateurs} (4 sem) - 30-50 utilisateurs
    \item \textbf{Finalisation} (3 sem) - Corrections, documentation
\end{enumerate}
\end{block}

\begin{center}
\textbf{Durée totale : 6 mois}
\end{center}
\end{frame}

% ============================================================================
% CONCLUSION
% ============================================================================

\section{Conclusion}

\begin{frame}{Récapitulatif}
\begin{block}{Système complet implémenté}
\begin{itemize}
    \item[\textcolor{successgreen}{\faCheckCircle}] Reconnaissance faciale opérationnelle
    \item[\textcolor{successgreen}{\faCheckCircle}] Chatbot intelligent avec RAG
    \item[\textcolor{successgreen}{\faCheckCircle}] Transcription et enregistrement audio
    \item[\textcolor{successgreen}{\faCheckCircle}] Gestion réservations automatique
    \item[\textcolor{successgreen}{\faCheckCircle}] Interface admin complète
    \item[\textcolor{successgreen}{\faCheckCircle}] Intégration Robot Pepper
    \item[\textcolor{successgreen}{\faCheckCircle}] Notifications email avec QR codes
\end{itemize}
\end{block}

\begin{block}{Objectifs}
Améliorer l'expérience client tout en validant l'acceptabilité des technologies innovantes
\end{block}
\end{frame}

\begin{frame}{Merci pour votre attention}
\begin{center}
\Large
\textbf{Questions ?}

\vspace{2cm}

\large
Équipe de Développement\\
Université d'Avignon

\vspace{1cm}

\today
\end{center}
\end{frame}

\end{document}

